\section{Related Work}
\label{sec:related}

\paragraph{WASM against native, one thread.} Haas et al.~\cite{haas2017}
introduced WebAssembly and reported it within 10\% of native on most of
PolyBenchC. Jangda et al.~\cite{jangda2019} ran the SPEC CPU suites in
Chrome and Firefox through a kernel extension that gives the browser a Unix
interface, and measured WASM 45\% slower than native on average in Firefox
and 55\% in Chrome; they traced it to more instructions, more loads and
stores, and more branches, from register pressure and safety checks. Yan et
al.~\cite{yan2021} showed that the compiler, its optimization flags, the
input size and the browser all change WASM results, sometimes in opposite
directions. Szewczyk et al.~\cite{szewczyk2022} measured the bounds checks
of linear memory outside browsers: under 20\% in runtimes that use virtual
memory, up to 650\% in others. All four are single-threaded. We take their
lesson for the design (pin the compiler, the flags, the input and the
browser version; report absolute ratios) and add threads. Our one-thread
ratios, 1.17 for Chrome and 1.36 for Firefox on the filter, are smaller than
Jangda's averages; we measure where they come from as instructions and
cycles per pixel (Section~\ref{sec:why-rq1}).

\paragraph{WASM multithreading.} Watt et al.~\cite{watt2019} gave the
formal semantics of WASM threads and atomics that this study relies on for
``always sequentially consistent''. S{\'a}nchez-Artigas and
Bohoyo-Bengoetxea~\cite{sanchez2025} studied portable WASM multithreading
outside the browser and found overheads in the allocator and in atomics;
this is why our allocation stays outside the timed region and why atomics
are the factor of RQ2. Closest to us, Kyriakou and
Tselikas~\cite{kyriakou2022} measured Rust compiled to WASM in Chromium at
93.5\% of a native Node.js addon single-threaded and 67.9\% multithreaded:
one workload without synchronization, 1 against 24 threads, means only. We
add native code as the baseline in both engines, a sweep over thread counts,
synchronization as a controlled factor, inferential statistics and a tested
cause.

\paragraph{Atomics and contention.} Schweizer et al.~\cite{schweizer2015}
measured atomic operations on several architectures and showed that on x86
compare-and-swap, fetch-and-add and swap cost about the same, whatever their
semantics, and that their cost under contention is that of moving the cache
line. Our RQ2 design takes its mechanism from them (padded counters, so that
what is measured is true sharing), and our finding qualifies their result:
the instructions around the atomic operation, not its kind, decide the cost
(Section~\ref{sec:why-rq2}).

\paragraph{Method.} Georges et al.~\cite{georges2007} showed that
conclusions about managed runtimes change with the statistics used and
asked for independent VM invocations and confidence intervals; Kalibera and
Jones~\cite{kalibera2013} added that repetition has levels (iterations
inside an invocation, invocations) and that the effort belongs at the level
where the variance is. Our unit of analysis is the run, a fresh process or
browser launch, with iterations summarized inside it. Mytkowicz et
al.~\cite{mytkowicz2009} showed that link order and environment size alone
move measurements; we therefore never mix runs of two builds and randomize
the order of all runs. Hoefler and Belli~\cite{hoefler2015} listed rules for
reporting parallel performance, among them: report the baseline of a
speedup, report absolute performance, use ratios and geometric means
consistently, and state the confidence of every number. The robustness check
uses the test of Mann and Whitney~\cite{mann1947}.
